
\documentclass{article}
\usepackage{colortbl}
\usepackage{makecell}
\usepackage{multirow}
\usepackage{supertabular}

\begin{document}

\newcounter{utterance}

\centering \large Interaction Transcript for game `eqbench', experiment `eq\_bench', episode 7 with reflexion{-}Qwen3.5{-}4B.
\vspace{24pt}

{ \footnotesize  \setcounter{utterance}{1}
\setlength{\tabcolsep}{0pt}
\begin{supertabular}{c@{$\;$}|p{.15\linewidth}@{}p{.15\linewidth}p{.15\linewidth}p{.15\linewidth}p{.15\linewidth}p{.15\linewidth}}
    \# & \multicolumn{2}{c}{Player} && \multicolumn{2}{c}{Game Master} \\
    \hline

    \theutterance \stepcounter{utterance}  
    & & & \multicolumn{4}{p{0.6\linewidth}}{
        \cellcolor[rgb]{0.9,0.9,0.9}{
            \makecell[{{p{\linewidth}}}]{
                \texttt{\tiny{[P1$\langle$GM]}}
                \texttt{You have attempted this same task before and did not succeed. These are the plans you wrote after those attempts.} \\
\texttt{Trial \#1: New plan: I failed because I assigned a score of 0 to "Relief," violating the rule that I must score at least one emotion > 0. Although I correctly identified Resentment and Frustration as the primary emotions, I need to ensure that at least one of the other options (Relief, Surprise, or Frustration) receives a non{-}zero score. I will re{-}evaluate the dialogue to see if there is any element of surprise or frustration that warrants a score greater than zero, or simply adjust the existing scores to ensure the constraint is met while maintaining the relative intensity of the emotions. For instance, I could give Frustration a score of 10 to be safe, or give Surprise a low score like 1 if the sudden shift in tone was unexpected. The key is to avoid a total zero across all categories.} \\
\\ 
\texttt{New plan: I will assign a score of 10 to Frustration and 8 to Resentment, and 0 to Relief and Surprise, ensuring that at least one emotion has a score greater than zero. I will double{-}check the output format to ensure no extra commentary is included.} \\
\\ 
\texttt{Trial \#2: Relief: 0} \\
\texttt{Resentment: 8} \\
\texttt{Surprise: 1} \\
\texttt{Frustration: 10} \\
\\ 
\texttt{Trial \#3: New plan: I will assign a score of 10 to Frustration and 8 to Resentment, and 0 to Relief and Surprise, ensuring that at least one emotion has a score greater than zero. I will double{-}check the output format to ensure no extra commentary is included.} \\
\\ 
\\ 
\\ 
\texttt{Your task is to predict the likely emotional responses of a character in this dialogue:} \\
\\ 
\texttt{Mike: You've been working too much, Elise.} \\
\texttt{Elise: And you've been reading too much. We all have our vices, Mike.} \\
\texttt{Mike: It's not a vice, it's a concern. You're burning yourself out.} \\
\texttt{Elise: Oh, so now you're my therapist? } \\
\texttt{Mike: No, I'm your friend. And as your friend, I'm worried about you.} \\
\texttt{Elise: Well, 'friend', I didn't ask for your worry.} \\
\texttt{{[}End dialogue{]}} \\
\\ 
\texttt{At the end of this dialogue, Mike would feel...} \\
\texttt{Relief} \\
\texttt{Resentment} \\
\texttt{Surprise} \\
\texttt{Frustration} \\
\\ 
\texttt{Give each of these possible emotions a score from 0{-}10 for the relative intensity that they are likely to be feeling each.} \\
\\ 
\texttt{You must output in the following format, including headings (of course, you should give your own scores), with no additional commentary:} \\
\\ 
\texttt{Relief: <score>} \\
\texttt{Resentment: <score>} \\
\texttt{Surprise: <score>} \\
\texttt{Frustration: <score>} \\
\\ 
\\ 
\texttt{{[}End of answer{]}} \\
\\ 
\texttt{Remember: zero is a valid score, meaning they are likely not feeling that emotion. You must score at least one emotion > 0.} \\
\\ 
\texttt{Your answer:} \\
            }
        }
    }
     \\ \\

    \theutterance \stepcounter{utterance}  
    & \multicolumn{4}{p{0.6\linewidth}}{
        \cellcolor[rgb]{0.561,0.737,0.561}{
            \makecell[{{p{\linewidth}}}]{
                \texttt{\tiny{[P1$\rangle$GM]}}
                \texttt{Relief: 0} \\
\texttt{Resentment: 8} \\
\texttt{Surprise: 0} \\
\texttt{Frustration: 10} \\
            }
        }
    }
    & & \\ \\

    \theutterance \stepcounter{utterance}  
    & & & \multicolumn{2}{p{0.3\linewidth}}{
        \cellcolor[rgb]{0.9,0.9,0.9}{
            \makecell[{{p{\linewidth}}}]{
                \texttt{\tiny{[GM$|$GM]}}
                \texttt{Relief: 0} \\
\texttt{Resentment: 8} \\
\texttt{Surprise: 0} \\
\texttt{Frustration: 10} \\
            }
        }
    }
    & & \\ \\

    \theutterance \stepcounter{utterance}  
    & & & \multicolumn{2}{p{0.3\linewidth}}{
        \cellcolor[rgb]{0.9,0.9,0.9}{
            \makecell[{{p{\linewidth}}}]{
                \texttt{\tiny{[GM$|$GM]}}
                \texttt{Relief: 0} \\
\texttt{Resentment: 3} \\
\texttt{Surprise: 2} \\
\texttt{Frustration: 8} \\
            }
        }
    }
    & & \\ \\

    \theutterance \stepcounter{utterance}  
    & & & \multicolumn{2}{p{0.3\linewidth}}{
        \cellcolor[rgb]{0.9,0.9,0.9}{
            \makecell[{{p{\linewidth}}}]{
                \texttt{\tiny{[GM$|$GM]}}
                \texttt{game\_result = LOSE} \\
            }
        }
    }
    & & \\ \\

    \theutterance \stepcounter{utterance}  
    & & \multicolumn{4}{p{0.6\linewidth}}{
        \cellcolor[rgb]{0.95,0.95,0.95}{
            \makecell[{{p{\linewidth}}}]{
                \texttt{\tiny{[P1$\leftrightarrow$P]}}
                \texttt{Post{-}episode Reflexion (not sent to GM):} \\
\texttt{Relief: 0} \\
\texttt{Resentment: 8} \\
\texttt{Surprise: 1} \\
\texttt{Frustration: 10} \\
            }
        }
    }
    & \\ \\

\end{supertabular}
}

\end{document}
